% TOP_PROBLEM: {"id": 1442, "title": "Teschner's Bondage Number Conjecture", "category_id": 3, "category": {"id": 3, "name": "graph_theory", "display_name": "Graph Theory", "description": "Problems involving graphs, networks, and their properties.", "slug": "graph-theory", "order_index": 3, "created_at": "2026-07-31T15:26:25.671Z"}, "status": "open"}
% FIRST_POSTED: null
% ENTRY_KIND: statement_only
% Imported from: batch07.tex; UnsolvedMath ID 1442
\documentclass[11pt]{article}
\usepackage[margin=25mm]{geometry}
\usepackage{fontspec}
\setmainfont{Latin Modern Roman}
\usepackage{amsmath,amssymb,mathtools}
\usepackage{unicode-math}
\setmathfont{Latin Modern Math}
\usepackage{xurl,hyperref}
\hypersetup{colorlinks=true,urlcolor=blue}
\setcounter{secnumdepth}{0}
\setlength{\parindent}{0pt}
\setlength{\parskip}{5pt}
\setlength{\emergencystretch}{3em}
\newcommand{\sref}[2]{\hyperlink{source:#1:#2}{[S#2]}}
\newcommand{\cataloguescope}[1]{\paragraph{Recorded target.}#1}
\begin{document}
% UnsolvedMath ID 1442; source catalogue code GRAPH-055
\setcounter{section}{1441}
\section{Teschner's Bondage Number Conjecture}\label{1442}
{\small\sffamily UnsolvedMath ID 1442 · Graph Theory}
\subsection{Definitions and mathematical statement}

\paragraph{Shared notation.}$\mathbb N=\{1,2,\ldots\}$, and $\mathbb Z,\mathbb Q,\mathbb R,\mathbb C$ denote the integers, rationals, reals and complex numbers. Any other conventions are specified locally.


Let $G=(V,E)$ be a finite simple undirected graph with at least one edge. A set $D\subseteq V$ dominates $G$ when every vertex outside $D$ is adjacent to a vertex in $D$. The domination number $\gamma(G)$ is the minimum size of a dominating set. For $F\subseteq E$, the graph $G-F$ has the same vertices and the remaining edges. The bondage number is
\[
 b(G)=\min\{|F|:F\subseteq E,\ \gamma(G-F)>\gamma(G)\}.
\]
This minimum exists because deleting every edge gives domination number $|V|$, while a graph with an edge has smaller domination number. Write $\Delta(G)$ for maximum vertex degree.
\paragraph{Statement recorded in the supplied source.}
Does every finite simple graph with at least one edge satisfy
\[
 b(G)\le\frac32\Delta(G),
\]
equivalently $b(G)\le\lfloor3\Delta(G)/2\rfloor$? The floor is the greatest integer not exceeding its argument. No connectedness, planarity or surface-embedding assumption is imposed. Edgeless graphs, on which no edge removal can increase domination number, are outside the domain of the bondage-number minimum. \sref{1442}{2}
\subsection{Short English statement}
Can the domination number of every graph with an edge be increased by removing at most three halves of its maximum degree in edges?
\subsection{Sources}
\begin{description}
\item[\hypertarget{source:1442:1}{S1}] ULAM.AI and the UnsolvedMath Contributors, UnsolvedMath: A Curated Collection of Open Mathematics Problems, record 1442 (GRAPH-055). CC BY 4.0. \url{https://huggingface.co/datasets/ulamai/UnsolvedMath}.
\item[\hypertarget{source:1442:2}{S2}] Andrei Gagarin and Vadim Zverovich, The Bondage Number of Graphs on Topological Surfaces and Teschner’s Conjecture, arXiv:1209.1362v1 (2012), Section 1, definitions and Conjecture 1, pages 1--2. \url{https://arxiv.org/pdf/1209.1362v1}.
\end{description}
\label{end:1442}
\end{document}
